\co{what did you explore. Try to be sequential. I did first a temperature dependent experiment
    since the results were outstanding I then performed another experiment for testing what I say in ch4 was just culo or a consequence of my novel system. Ultimately, I have put at work one longstanding issue in electronics, i.e. is it possible to engineer the band energy such that I can obtain a p--n junction? How did you do that, I engineered the energy band in silico first. I tested the effect is minimum and that possibly a more viable route is the introduction in the lattice of a dopant atom.}

\co{start by elucidating your first experiment. What have you used and why.}



% Since it would be desirable to have an electronic device operating at room temperature, I then first studied a novel class of molecular materials, namely \glspl{MGNR}, with which I realised the first room-temperature-operating MGNR-based \gls{SET}, combined with the comprehension of a few aspects of the electronic transport through molecularly-made graphene nanoribbons never reported before.

% Since it would be desirable to have an electronic device operating at room
% temperature, I then first studied a novel class of molecular materials, namely
% MGNRs. Thanks to a novel approach which uses sterically-hindred molecular
% functionalisations grafted onto the edges, we have been able to successfully realised
% the first room-temperature-operating GNRSET, combined with the comprehension
% of a few aspects of the electronic transport through molecularly-made graphene
% nanoribbons, never reported before.


% In this thesis, I have investigated the electronic transport properties at the device level of molecular graphene nanoribbons, focusing on those that are relevant to their potential use as a quantum technology platform and as channel material in miniaturised all-carbon logic components. In particular, I sought to verify the existance electronic states robust enough to overcome the thermal broadening and to address the possibility to tune the electronic properties by adding chemical groups either along the edges or within the honeycomb lattice.

% We established in Ch.~\ref{ch:gnrset} that the temperature dependence of graphene nanoribbon conductance follows a power-law, and that the molecular precision with which \glspl{MGNR} can currently be synthesised allows quantum effects such as tunnelling events to be detected at temperatures close to room temperature. We attributed this behaviour to the length/width ratio of the ribbon strands, as well as the fact that long-range finite Coulomb interactions appear to play a role when measuring electronic properties, which was previously unreachable with STM experiments. A further step would be to determine whether the observed behaviour for \glspl{MGNR} is due to phase transitions that they undergo to or to the presence of defects inside the MGNR lattice. To accomplish this, a longer ribbon would need to be studied, as well as possibly adding another gate at the bottom to see if it is possible to create a double dot in a molecular system for the first time, which would represent a significant step forward in the creation of molecular qubits operating at the device level.

% We demonstrated in Ch.~\ref{ch:solubility} that the insertion of 3D bulky groups along the edges could be a helpful method for affording MGNR-based electrical devices where electronic and vibrational states can be examined, thus increasing the reproducibility of these types of devices. We were able to observe a strong suppression of the tunnelling current at low bias in combination with vibron-assisted excitations using milliKelvin transport spectroscopy, which was explained within the Franck--Condon model supported by \gls{DFT} studies about the changes induced in the electronic and phonon dispersions by the maleimide groups. Our findings indicate that the vibrational modes in \gls{GNR} are more tightly related to the electric charge than was previously reported for a class of devices built using suspended carbon nanotubes. Our statistical examination of the role of the bulky groups on a variety of devices was inconclusive due to the small sample size and the limited number of devices that could be cooled at the same time. It would thus be necessary to perform a more thorough statistical analysis using a cryogenic multiplexer, which would allow us to increase the number of measurable devices with a single cooldown.
% From a basic and technological standpoint, the discoveries gained in this chapter pave the way for the employment in spintronic devices not only of localised magnetic centres, as in single molecule magnets, but also of magnetic states of delocalised electronic systems, such as graphene nanoribbons.

% In Ch.~\ref{gnr-hybrid:sec:intro}, we looked at two distinct methods for doping graphene nanoribbons. We predict that if large changes in the Fermi energy position of the graphene nanoribbon energy level are desired, chemical groups acting on the $\pi$-electron system should be used, and that electron donating groups induce relative position adjustments that will be extremely difficult to detect at the device level. However, the incorporation of porphyrin complexes within the graphene honeycomb appears to be a more valuable choice for coupling the spin and charge degrees of freedom without losing the electrical properties. We discover that the intrinsic mobilities caused by the addition of Ni-based porphyrin rings are comparable to those reported for pure MGNR, that the switching behaviour for certain devices is low enough to use these hybrids as channel material in FET devices, and that the off-state currents are already competitive with those of most CNT-based FETs.

% Finally, we note three issues and open concerns that emerged from this research and should be addressed in future ones. All the graphene nanoribbons used in this thesis were synthesised with methods producing ribbons of different lengths. This can be issue because the reproducibility for both \glspl{SET} and \glspl{FET} depends on the lateral extension of the \glspl{MGNR}. Moreoverm, longer ribbon would be required for the construction of microprocessors, as as the one recently created entirely from carbon nanotubes~\cite{hills2019modern}. Studies exploring ribbons from a particular separation chromatography could elucidate and confirm the presence or absence of finite-range interaction and whether there is a cutoff length to be fulfilled in order to verify if length is a critical parameter. Second, while the formation of GNR--GNR heterojunctions has demonstrated the formation of topologically non-trivial phases~\cite{lee2018topological}, it would be interesting to see if the introduction of chemical groups along the edges or inside the honeycomb lattice will induce a particle--hole symmetry breaking, and in the formation of new topologically non-trivial phases. Finally, the ability to perform quantum operations on single molecular qubits at the device level remains elusive. Graphene nanoribbons could be employed in a single-molecule EPR transport experiment, where microwave pulses excite and control the electron-spin coherence, and whether topologically-protected electronic states neare the Fermi energy could be used to improve their resistance agains main source or dechorence~\cite{lombardi2019quantum}. Furthermore, it would be interesting to insert graphene nanoribbons between two superconducting leads to see if the absence of backscattering inherent in graphene-derived nanostructures could help with Cooper-pair transmission.

% MGNRs host a plethora of physics phenomena which they should be first investigate thouroghly, both through STM and at the device level. The necessity of synthesising longer ribbon while mantaining a high degree of debundling will also be fundamental towards the realisation of a first MGNR-based microprocessor.
